\centerline{\bf \huge 8 класс}
\centerline{\normalsize{\it Работа рассчитана на \bf{180} минут}}

\problem{8.1}
В десятизначном числе пронумеровали слева направо все цифры числами от 1 до 10 , после чего все цифры с чётными номерами в том же порядке переставили на первые пять мест, а все цифры с нечётными номерами в том же порядке переставили на последние 5 мест. Получившееся число совпало с исходным. Какое наибольшее количество различных цифр может быть в таком числе?

\problem{8.2}
На биссектрисе острого угла $A B C$ отмечены точки $D$ и $E$ так, что $A B=B D=A E$ и $D E=B C$ (точка $D$ лежит между $B$ и $E$ ). Докажите, что $A D=C D$.

\problem{8.3}
На доске написаны четыре числа: $2,3,4$ и 6 . За один шаг можно выбрать любые три из них, первое умножить на 2 , второе - на 3 , а третье - на 6 (при этом три старых числа стирают, а на их место записывают три новых). Можно ли через несколько шагов получить на доске четыре равных числа?

\problem{8.4}
Найдите значение суммы

$$
\begin{gathered}
S=\left(1^2+1 \cdot 2+2^2\right)+\left(2^2+2 \cdot 3+3^2\right)+\left(3^2+3 \cdot 4+4^2\right)+\cdots \\
+\left(98^2+98 \cdot 99+99^2\right)+\left(99^2+99 \cdot 100+100^2\right)
\end{gathered}
$$

\problem{8.5}
На столе лежат 100 карточек с числами от 1 до 100 . Двое играют в следующую игру. Ходят по очереди. За один ход можно взять со стола любую карточку. Игра заканчивается, когда на столе останется две карточки. Второй выигрывает, если числа на оставшихся карточках отличаются ровно на 10 . Иначе выигрывает первый. Кто выигрывает при правильной игре?